\chapter{Methodology}

\section{System Architecture}
The NS-CausalKT model is designed as a three-layered architecture. Each layer addresses a specific aspect of the learning process: neural sequence modeling, symbolic constraint enforcement, and causal intervention.

\subsection{Neural Student Model (NSM)}
The NSM uses a 2-layer Transformer encoder with multi-head attention. For each student interaction $(q_t, a_t)$, we create an embedding $e_t$:
\begin{equation}
    e_t = \text{Embed}(q_t) \oplus \text{Embed}(a_t) \oplus \text{Embed}(c_t) \oplus \Delta t_t
\end{equation}
The Transformer output $H$ is then refined through a GRU to capture temporal dynamics like the spacing effect (forgetting curves).

\subsection{Symbolic Concept Layer (SCL)}
The SCL ensures that the latent knowledge state $\mu_t$ follows the rules of the domain. We define a Directed Acyclic Graph (DAG) $G$ representing concept prerequisites. We use **Logic Tensor Networks (LTN)** to define soft rules. For example, the "Prerequisite Necessity" rule states:
\begin{equation}
    \forall s, t: \text{masters}(s, c_i, t) \implies \text{masters}(s, c_j, t) \text{ if } (c_j \to c_i) \in E
\end{equation}
This rule is enforced using the Łukasiewicz t-norm, making the logic differentiable and suitable for gradient-based training.

\subsection{Causal SCM Layer (CSL)}
The CSL implements a Structural Causal Model (SCM). The Average Causal Effect (ACE) of a concept $c_j$ on the student's performance $a_{t+1}$ is computed using Pearl's do-calculus:
\begin{equation}
    \text{ACE}(c_j \to a_{t+1}) = E[a_{t+1} | do(c_j=1)] - E[a_{t+1} | do(c_j=0)]
\end{equation}
By computing ACE for all prerequisite concepts, we can identify the "Root Cause" of a student's failure.

\section{Loss Functions}
The model is trained end-to-end using a joint objective function $L_{total}$:
\begin{equation}
    L_{total} = L_{pred} + \lambda_1 L_{sym} + \lambda_2 L_{causal} + \lambda_3 L_{smooth}
\end{equation}
Where:
\begin{itemize}
    \item $L_{pred}$ is the Binary Cross-Entropy loss for performance prediction.
    \item $L_{sym}$ ensures consistency with the prerequisite graph.
    \item $L_{causal}$ regularizes the mastery state based on causal constraints.
    \item $L_{smooth}$ penalizes sudden, unrealistic jumps in knowledge mastery.
\end{itemize}

\section{Data Preprocessing}
The ASSISTments 2009 dataset was preprocessed to remove noise and ensure consistency.
\begin{itemize}
    \item \textbf{Sequence Length}: Sequences were truncated or padded to a fixed length of 200 interactions.
    \item \textbf{Concept Mapping}: Each problem was mapped to a unique skill ID.
    \item \textbf{Splitting}: A student-level 80/20 split was used for training and validation to prevent data leakage.
\end{itemize}
